\documentclass[12pt,a4paper]{article}

\usepackage[a4paper,margin=1in]{geometry}
\usepackage{setspace}
\usepackage{enumitem}
\usepackage{hyperref}
\usepackage{titlesec}
\usepackage{fancyhdr}

% Page settings
\setstretch{1.15}
\setlength{\parindent}{0pt}
\setlength{\parskip}{6pt}

% Header and Footer
\pagestyle{fancy}
\fancyhf{}
\lhead{Tourism Guide Planner}
\rhead{UCS503}
\cfoot{\thepage}

% Section formatting
\titleformat{\section}
{\large\bfseries}
{\thesection.}
{0.5em}
{}

\titleformat{\subsection}
{\normalsize\bfseries}
{\thesubsection}
{0.5em}
{}

% Hyperlinks
\hypersetup{
    colorlinks=true,
    linkcolor=blue,
    urlcolor=blue
}

\begin{document}

% =========================================================
% TITLE PAGE
% =========================================================

\begin{titlepage}

\begin{center}

\vspace*{0.7cm}

{\Huge \textbf{Tourism Guide Planner}}

\vspace{0.6cm}

{\Large \textbf{Project Proposal for UCS503}}

\vspace{1cm}

{\large Submitted to:}

\vspace{0.2cm}

{\Large \textbf{Ms. Paramveer Kaur}}

\vspace{0.4cm}

{\Large \textbf{Thapar Institute of Engineering and Technology}}

\vspace{1.2cm}

\begin{tabular}{ccc}

\textbf{Tanish Bansal}
&
\textbf{Ayush Jindal}
&
\textbf{Mehak Garg}
\\[0.25cm]

Roll No.: 1024030742
&
Roll No.: 1024030748
&
Roll No.: 1024030739
\\[0.15cm]

tbansal1\_be24@thapar.edu
&
ajindal1\_be24@thapar.edu
&
mgarg3\_be24@thapar.edu

\end{tabular}

\vfill

{\large \textbf{Department of Computer Science and Engineering}}

\vspace{0.25cm}

{\large Thapar Institute of Engineering and Technology}

\vspace{0.25cm}

{\large 2026}

\end{center}

\end{titlepage}


% =========================================================
% TABLE OF CONTENTS
% =========================================================

\tableofcontents

\newpage


% =========================================================
% 1. INTRODUCTION
% =========================================================

\section{Introduction}

Tourism Guide Planner is a web-based platform designed to help foreign
tourists explore historical places in an easy and informative way.

Tourists often need to visit different websites to find information about
historical places, understand their importance, find local guides, and
plan their visit.

Our project aims to provide these services through a single platform.
Users will be able to search for historical places, read information about
them, use a digital guide book, interact with an AI chatbot, and find
local guides.

The platform will also allow tourists to upload photographs from their
visits, making the experience more interactive.


% =========================================================
% 2. PROBLEM STATEMENT
% =========================================================

\section{Problem Statement}

Foreign tourists visiting historical places may face several difficulties
while planning and experiencing their trip.

The major problems are:

\begin{itemize}[leftmargin=*]

    \item Historical information is available on different websites.

    \item Tourists may find it difficult to understand the historical
    importance of a place.

    \item Finding a reliable local guide can be difficult.

    \item Tourists may not have immediate help when they have questions
    about a historical place.

    \item Information about guides, historical places, and tourism services
    is not available in one place.

\end{itemize}

Tourism Guide Planner aims to solve these problems by providing a
centralized platform for historical tourism.


% =========================================================
% 3. OBJECTIVES
% =========================================================

\section{Objectives}

The main objectives of the project are:

\begin{itemize}[leftmargin=*]

    \item To provide information about historical places.

    \item To allow tourists to search and explore historical destinations.

    \item To provide a digital guide book.

    \item To provide an AI chatbot for historical questions.

    \item To help tourists find and hire local guides.

    \item To provide guide information such as experience, languages,
    ratings, and reviews.

    \item To allow tourists to upload photographs.

    \item To provide all major tourism features through one platform.

\end{itemize}


% =========================================================
% 4. PROPOSED SYSTEM
% =========================================================

\section{Proposed System}

The proposed system is a web application that combines historical
information and tourism services.

A tourist can first search for a historical place and view its details.
The tourist can then use the digital guide book to learn more about the
location.

If the tourist needs additional information, the AI chatbot can answer
questions related to the historical place.

The tourist can also search for local guides and view their profiles
before sending a booking request.

After visiting the location, the tourist can upload photographs and share
their experience.


% =========================================================
% 5. MAIN FEATURES
% =========================================================

\section{Main Features}

\subsection{Historical Places}

The platform will provide information about different historical places.

The information can include:

\begin{itemize}[leftmargin=*]

    \item Name
    \item Location
    \item Historical background
    \item Cultural importance
    \item Architecture
    \item Important events
    \item Interesting facts
    \item Images

\end{itemize}


\subsection{Digital Guide Book}

The digital guide book will provide detailed information about a
historical place.

Tourists can read about the history, architecture, culture, and important
events associated with the location.


\subsection{AI Chatbot}

The AI chatbot will allow tourists to ask questions about historical
places.

For example:

\begin{itemize}[leftmargin=*]

    \item Who built this monument?
    \item When was it constructed?
    \item Why is it important?
    \item What happened here?
    \item What should I see at this place?

\end{itemize}

The chatbot will provide simple and understandable answers.


\subsection{Local Guide Search}

Tourists will be able to search for local guides.

Guide profiles can contain:

\begin{itemize}[leftmargin=*]

    \item Name
    \item Experience
    \item Languages spoken
    \item Rating
    \item Reviews
    \item Availability

\end{itemize}


\subsection{Guide Booking}

After selecting a guide, tourists can choose a suitable date and time
and submit a booking request.

The guide can then accept or reject the request.


\subsection{Photo Upload}

Tourists will be able to upload photographs taken during their visit.

This feature will make the platform more interactive and allow tourists
to share their experiences.


% =========================================================
% 6. BASIC WORKFLOW
% =========================================================

\section{Basic System Workflow}

The basic working of the system will be:

\begin{enumerate}[leftmargin=*]

    \item Tourist creates an account or logs into the platform.

    \item Tourist searches for a historical place.

    \item Tourist views information about the place.

    \item Tourist can read the digital guide book.

    \item Tourist can ask questions using the AI chatbot.

    \item Tourist can search for local guides.

    \item Tourist selects a guide and sends a booking request.

    \item Guide accepts or rejects the booking.

    \item Tourist visits the historical place.

    \item Tourist can upload photographs after the visit.

\end{enumerate}


% =========================================================
% 7. TECHNOLOGY STACK
% =========================================================

\section{Technology Stack}

The following technologies can be used to develop the project.

\subsection{Frontend}

\begin{itemize}[leftmargin=*]

    \item React.js / Next.js
    \item HTML
    \item CSS
    \item JavaScript
    \item Tailwind CSS

\end{itemize}


\subsection{Backend}

\begin{itemize}[leftmargin=*]

    \item Node.js
    \item Express.js
    \item REST APIs

\end{itemize}


\subsection{Database}

A database such as MongoDB or PostgreSQL can be used to store:

\begin{itemize}[leftmargin=*]

    \item Users
    \item Historical places
    \item Guides
    \item Bookings
    \item Reviews
    \item Photographs

\end{itemize}


\subsection{Artificial Intelligence}

An AI/LLM API can be used to develop the historical chatbot.

The chatbot can use historical information stored in the system to
provide relevant answers.


\subsection{Version Control}

Git and GitHub will be used for project development and collaboration.

The project repository will contain the frontend, backend, database,
AI-related files, and documentation.


% =========================================================
% 8. USER ROLES
% =========================================================

\section{User Roles}

There will be three main types of users in the system.


\subsection{Tourist}

A tourist can:

\begin{itemize}[leftmargin=*]

    \item Register and log in.
    \item Search historical places.
    \item Read historical information.
    \item Use the digital guide book.
    \item Ask questions using the AI chatbot.
    \item Search for guides.
    \item Book guides.
    \item Upload photographs.

\end{itemize}


\subsection{Local Guide}

A local guide can:

\begin{itemize}[leftmargin=*]

    \item Create a profile.
    \item Add experience and languages.
    \item Add availability.
    \item Receive booking requests.
    \item Accept or reject bookings.
    \item Receive ratings and reviews.

\end{itemize}


\subsection{Administrator}

The administrator can:

\begin{itemize}[leftmargin=*]

    \item Manage users.
    \item Manage historical places.
    \item Manage guide profiles.
    \item Manage bookings.
    \item Manage reviews and uploaded content.

\end{itemize}


% =========================================================
% 9. FUTURE SCOPE
% =========================================================

\section{Future Scope}

The project can be expanded with several additional features.

\subsection{Audio Guide}

The digital guide book can be converted into an audio guide so that
tourists can listen to historical information while visiting a location.


\subsection{AI Voice Assistant}

The chatbot can be extended to support voice-based questions and answers.


\subsection{Multiple Languages}

The platform can support multiple languages to make it more useful for
foreign tourists.


\subsection{Mobile Application}

A mobile application can be developed for Android and iOS.

The mobile application can provide:

\begin{itemize}[leftmargin=*]

    \item GPS-based location services
    \item AI chatbot
    \item Audio guide
    \item Guide booking
    \item Photo uploads

\end{itemize}


\subsection{More Historical Destinations}

The system can initially contain selected historical places and can later
be expanded to cover more cities, states, and countries.


% =========================================================
% 10. TEAM RESPONSIBILITIES
% =========================================================

\section{Team Responsibilities}

The project consists of three members.

\begin{center}

\begin{tabular}{|p{4cm}|p{8cm}|}
\hline
\textbf{Team Member} & \textbf{Responsibility} \\
\hline

Tanish Bansal &
Frontend development and historical places module
\\
\hline

Ayush Jindal &
Backend development, database and guide booking
\\
\hline

Mehak Garg &
AI chatbot, digital guide book and photo upload
\\
\hline

\end{tabular}

\end{center}


% =========================================================
% 11. EXPECTED OUTCOME
% =========================================================

\section{Expected Outcome}

The expected outcome of the project is a working tourism website where
foreign tourists can:

\begin{itemize}[leftmargin=*]

    \item Discover historical places.

    \item Read historical information.

    \item Use a digital guide book.

    \item Ask questions using an AI chatbot.

    \item Find and book local guides.

    \item Upload photographs.

\end{itemize}

The system will provide a simple and centralized tourism experience.


% =========================================================
% 12. CONCLUSION
% =========================================================

\section{Conclusion}

Tourism Guide Planner aims to provide a simple and useful platform for
foreign tourists visiting historical places.

The system combines historical information, a digital guide book, an AI
chatbot, local guide booking, and tourist photo uploads.

Instead of using multiple platforms, tourists can access these features
from one website.

The project can also be extended in the future with audio guides,
multiple languages, AI voice assistance, and mobile applications.

Overall, Tourism Guide Planner will make the process of exploring
historical places more convenient, informative, and interactive.

\vspace{1cm}

\begin{center}

\textbf{Tourism Guide Planner}

\vspace{0.3cm}

\textit{Explore History. Discover Places. Connect with Local Guides.}

\end{center}

\end{document}